% Options for packages loaded elsewhere
\PassOptionsToPackage{unicode}{hyperref}
\PassOptionsToPackage{hyphens}{url}
\documentclass[
]{article}
\usepackage{xcolor}
\usepackage[margin=2cm]{geometry}
\usepackage{amsmath,amssymb}
\setcounter{secnumdepth}{-\maxdimen} % remove section numbering
\usepackage{iftex}
\ifPDFTeX
  \usepackage[T1]{fontenc}
  \usepackage[utf8]{inputenc}
  \usepackage{textcomp} % provide euro and other symbols
\else % if luatex or xetex
  \usepackage{unicode-math} % this also loads fontspec
  \defaultfontfeatures{Scale=MatchLowercase}
  \defaultfontfeatures[\rmfamily]{Ligatures=TeX,Scale=1}
\fi
\usepackage{lmodern}
\ifPDFTeX\else
  % xetex/luatex font selection
\fi
% Use upquote if available, for straight quotes in verbatim environments
\IfFileExists{upquote.sty}{\usepackage{upquote}}{}
\IfFileExists{microtype.sty}{% use microtype if available
  \usepackage[]{microtype}
  \UseMicrotypeSet[protrusion]{basicmath} % disable protrusion for tt fonts
}{}
\makeatletter
\@ifundefined{KOMAClassName}{% if non-KOMA class
  \IfFileExists{parskip.sty}{%
    \usepackage{parskip}
  }{% else
    \setlength{\parindent}{0pt}
    \setlength{\parskip}{6pt plus 2pt minus 1pt}}
}{% if KOMA class
  \KOMAoptions{parskip=half}}
\makeatother
\usepackage{longtable,booktabs,array}
\newcounter{none} % for unnumbered tables
\usepackage{calc} % for calculating minipage widths
% Correct order of tables after \paragraph or \subparagraph
\usepackage{etoolbox}
\makeatletter
\patchcmd\longtable{\par}{\if@noskipsec\mbox{}\fi\par}{}{}
\makeatother
% Allow footnotes in longtable head/foot
\IfFileExists{footnotehyper.sty}{\usepackage{footnotehyper}}{\usepackage{footnote}}
\makesavenoteenv{longtable}
\usepackage{graphicx}
\makeatletter
\newsavebox\pandoc@box
\newcommand*\pandocbounded[1]{% scales image to fit in text height/width
  \sbox\pandoc@box{#1}%
  \Gscale@div\@tempa{\textheight}{\dimexpr\ht\pandoc@box+\dp\pandoc@box\relax}%
  \Gscale@div\@tempb{\linewidth}{\wd\pandoc@box}%
  \ifdim\@tempb\p@<\@tempa\p@\let\@tempa\@tempb\fi% select the smaller of both
  \ifdim\@tempa\p@<\p@\scalebox{\@tempa}{\usebox\pandoc@box}%
  \else\usebox{\pandoc@box}%
  \fi%
}
% Set default figure placement to htbp
\def\fps@figure{htbp}
\makeatother
\setlength{\emergencystretch}{3em} % prevent overfull lines
\providecommand{\tightlist}{%
  \setlength{\itemsep}{0pt}\setlength{\parskip}{0pt}}
\usepackage{bookmark}
\IfFileExists{xurl.sty}{\usepackage{xurl}}{} % add URL line breaks if available
\urlstyle{same}
\hypersetup{
  hidelinks,
  pdfcreator={LaTeX via pandoc}}

\author{}
\date{}

\begin{document}

\section{Does a boosted model beat the best simple baseline?
Endpoint-matched evaluation of daily retail demand forecasts on public
data}\label{does-a-boosted-model-beat-the-best-simple-baseline-endpoint-matched-evaluation-of-daily-retail-demand-forecasts-on-public-data}

\emph{Working title. Draft for author review --- not submitted. Author
names, affiliations and declarations: {[}AUTHOR TO COMPLETE{]}.}

\subsection{Abstract}\label{abstract}

\textbf{Background.} A previous study in this project found that a
LightGBM model with a Tweedie objective did not beat a trailing 28-day
median on pooled item-day weighted absolute percentage error (WAPE) for
daily retail demand, while winning after demand was aggregated.
Forecasting theory attributes such reversals to the scoring function
(absolute error rewards the conditional median, squared error the mean),
but the size of the effect, and whether the earlier negative result
survives a candidate set that includes median-type objectives, had not
been measured under a stricter protocol. \textbf{Methods.} Direct 7- and
28-day-ahead forecasts of daily units were evaluated on three M5 stores
(350/350/349 department-stratified series) and the 100 largest-revenue
SKUs of Online Retail II, in six rolling 28-day test blocks after three
validation blocks, on a confirmatory window that no experiment had
scored (W2) and on replication windows (W3, W1). Eight boosted-tree or
ridge candidates and 31 simple baselines (16 estimators × matched and
fresh information) were compared. Selection was made on validation data
for the endpoint being reported (item-day or subset-day WAPE); the
protocol was hash-frozen before any confirmatory outcome was read.
Inference used a paired crossed bootstrap over series and weeks
(secondary: two-way cluster-robust interval), Holm correction, a 2 \%
minimum improvement and a dual-method rule. \textbf{Results.} On the
confirmatory window the validation-selected absolute-error LightGBM beat
the best same-information baseline on item-day WAPE in all 8
dataset--horizon cells (median +4.3 \%, range +2.5 \% to +6.0 \%), and
the improvement was positive in all 22 cells over three windows
(supported in 20). It beat the earlier Tweedie model in every cell (+7.9
\% to +16.8 \%). Against baselines that see up to 28 days more recent
data the advantage was not established at the 28-day horizon. On the
aggregate (subset-day) endpoint the selected models were mean-type
(Tweedie or squared error); the item-day winner, with a relative bias of
−0.29 to −0.39, lost to simple baselines (subset-day WAPE 0.35 against
0.11 for CA\_1, 28 days), and the improvement of the aggregate-selected
model over the best baseline was supported in 4 of 8 cells (point
estimates positive in 7 of 8; one boundary cell was re-examined with a
supplementary calculation that left every verdict unchanged).
\textbf{Conclusions.} For these samples the earlier negative result
changes sign when the objective and the evaluation endpoint are matched,
with modest gains at item-day and endpoint-dependent rankings. The
finding is an empirical quantification of established theory on four
public data sets, not a new method. Statistical (ETS, ARIMA, Theta),
neural and probabilistic competitors were not evaluated, and the result
does not extend to other retailers or to decision costs. The
verification was internal (one AI assistant wrote the study and its
checks).

\textbf{Keywords:} demand forecasting; intermittent demand; gradient
boosting; scoring functions; forecast evaluation; aggregation;
reproducibility.

\subsection{1. Introduction}\label{introduction}

Boosted trees are strong forecasters of retail demand: in the M5
accuracy competition most of the examined methods, and almost all of the
top-50 methods for which a description was available, were
LightGBM-based {[}Makridakis et al., 2022{]}, and tree-based methods
outperformed other approaches there {[}Januschowski et al., 2022{]};
reviews of Kaggle forecasting competitions, whose data are typically
intermittent, report that gradient-boosted trees and neural networks are
strong {[}Bojer \& Meldgaard, 2021{]}. Yet strong aggregate results do
not settle whether a learned model beats a \emph{well-chosen simple
baseline} on a given error measure. Forecasting theory is explicit that
the optimal point forecast depends on the scoring function: squared
error is minimised by the conditional mean, absolute error by the
conditional median {[}Gneiting, 2011; Kolassa, 2020; Koenker \& Bassett,
1978{]}. For sparse daily demand, whose conditional median is often
zero, absolute-error measures therefore reward low forecasts (a
consequence of the result just cited, not a finding of this study);
conventional accuracy measures can also rank forecasters of intermittent
demand incorrectly or be degenerate {[}Prestwich et al., 2014; Hyndman
\& Koehler, 2006{]}, and the evaluation and optimisation of
intermittent-demand forecasts has its own literature {[}Kolassa, 2016;
Kourentzes, 2014{]}. Aggregating demand over time or items changes the
forecasting problem {[}Nikolopoulos et al., 2011; Athanasopoulos et al.,
2017{]}.

A previous study in this project (the ``v1 study'') found that a
LightGBM model with a Tweedie objective did not beat a trailing 28-day
median in pooled item-day WAPE on two M5 stores and one retail dataset,
while winning after demand was summed over days. An internal audit of
that study (carried out by the same AI assistant, so not independent)
showed that the contrast was largely the baseline's level bias, that no
candidate model targeted a median, and that the evaluation windows had
been examined repeatedly. The present study asks the question again
under a stricter design.

\textbf{Research question.} Under model selection fixed on validation
data \emph{for the endpoint being reported}, and on evaluation windows
that no experiment had scored, does a boosted model beat the best simple
baseline; by how much, with what uncertainty; does the answer change
between the item-day and the aggregated endpoint; does the v1 negative
result survive when median-type objectives are in the candidate set; and
does the answer replicate across stores, horizons and windows?

\subsection{2. Related work and what is already
known}\label{related-work-and-what-is-already-known}

The relationship between loss function and optimal forecast {[}Gneiting,
2011; Kolassa, 2020{]}, the unreliability of conventional error measures
for intermittent and low-count demand {[}Hyndman \& Koehler, 2006;
Prestwich et al., 2014; Kolassa, 2016; Kourentzes, 2014{]}, temporal and
cross-sectional aggregation {[}Nikolopoulos et al., 2011; Athanasopoulos
et al., 2017; Hyndman et al., 2011; Spiliotis et al., 2021{]}, and the
strength of boosted trees on M5-type data {[}Makridakis et al., 2022;
Januschowski et al., 2022; Bojer \& Meldgaard, 2021; Spiliotis et al.,
2022{]} are established. Evidence that machine-learning methods do not
automatically beat statistical ones also exists {[}Makridakis et al.,
2018; Makridakis et al., 2020{]}, and retail forecasting has its own
practice and measurement issues {[}Fildes et al., 2022{]}. Evaluation
guidance is given by Hewamalage et al.~{[}2023{]}. A preprint compares
local and global models for intermittent series {[}Damato et al.,
2026{]}.

In the sources read for this study (a bounded literature review, not a
systematic one), I did not find a report that combines endpoint-matched
validation selection, never-scored windows chosen under a written
freeze, a large prespecified baseline set, and paired two-way
dependence-aware intervals for a boosted model on public retail data.
That is an absence in what was read, not proof of absence. The expected
direction of the main effect follows from theory; the study measures its
size and location.

\subsection{3. Contributions and scope}\label{contributions-and-scope}

\emph{Contributions (incremental and empirical).} (i) Magnitudes of the
item-day and aggregate differences between validation-selected boosted
models and the best of 31 simple baselines on four public datasets, two
horizons and replication windows, under a frozen protocol; (ii) a test
of whether the v1 negative result survives a candidate set that includes
median-type objectives; (iii) a decomposition of the baseline's level
bias; (iv) a pre-freeze simulation of the interval methods and a
dual-method decision rule; (v) a documented package with frozen
protocol, pinned environment and deterministic settings (a
fresh-environment reproduction has not been performed). \emph{Not
contributions:} a new model, metric or principle; any claim about Indian
SMEs or other populations; causal effects; customer churn; predictive
intervals.

\subsection{4. Data and data rights}\label{data-and-data-rights}

M5 {[}Makridakis et al., 2022; Kaggle competition data{]}: daily unit
sales for 350 (CA\_1, TX\_1) and 349 (WI\_1) department-stratified
random series per store, drawn among series whose first sale falls on or
before day d\_1100 (products launched later are excluded; seed
20261015). Online Retail II {[}Chen, 2012; CC BY 4.0{]}: the 100 SKUs
with the largest revenue before the first validation block of each
window, daily gross units after documented cleaning (cancellations,
non-positive quantity or price, non-product codes and 33,666 exact
duplicate lines removed; returns excluded). M5 items are evaluated on
days with a listed selling price (days without a listing carry zero
demand by construction; verified); this eligibility rule is an
assumption about ex-ante knowledge. \textbf{Rights:} the M5 competition
rules (read 2026-10-22) permit use of the Competition Data ``for
non-commercial purposes only, including \ldots{} for academic research
and education'' and prohibit transmitting, duplicating, publishing or
redistributing it to anyone not participating in the competition; this
study uses it for academic research only, and neither the data nor any
array that contains M5 sales (stored predictions include the actual
values) is released (Section 13). Whether forecasts alone may be shared
is not stated in the rules and is treated as not permitted unless the
rights holder agrees.

\subsection{5. Windows and protocol}\label{windows-and-protocol}

A window is defined by the number of trailing days removed \emph{before
anything is built}: W1 (0), W2 (252), W3 (504). On the truncated series
the last 168 days form six consecutive 28-day test blocks and the
preceding 84 days three validation blocks. \textbf{W2 (all four
datasets) is the confirmatory window}; W3 (M5) and W1 for the fresh
store WI\_1 are replications; W1 for CA\_1, TX\_1 and Online Retail II
is a non-independent re-evaluation (those periods were scored
repeatedly). Exposure of every window to earlier experiments, including
a retired benchmark that used the W2/W3 \emph{periods} on 40 other M5
series, is documented (\texttt{WINDOW\_USE\_AUDIT.md}); no evaluated
entity's W2/W3 target was scored before. The design was fixed with
knowledge of the v1 results and the audit.

\emph{Forecasting task.} Direct forecasts of daily units for a minimum
lag H ∈ \{28, 7\} days: target day t uses sales history up to t − H plus
calendar and (M5) listed-price information of day t; a global model is
refit for each block on rows whose target day lies in {[}o − 730, o −
1{]} (o = block start). This is \textbf{not} the M5 competition setting.
Predictions are clipped at zero.

\emph{Freeze.} The protocol, code and tests were hashed on 2026-10-22
(08:15:19Z) before any W2/W3/WI\_1 validation or test outcome was read
(\texttt{FREEZE.json}); this is an \textbf{internal} freeze, not a
public pre-registration. Pre-freeze activity (unit tests, a dry run on
already-scored Online Retail II W1 data, a simulation on synthetic data,
literature review) is logged in \texttt{DEVIATIONS.md}.

\subsection{6. Models and baselines}\label{models-and-baselines}

\emph{Candidates} (shared base hyper-parameters: 500 trees (XGBoost
400), learning rate 0.05, 31 leaves, subsample 0.8, column sample 0.8):
LightGBM {[}Ke et al., 2017; gradient boosting: Friedman, 2001{]} with
Tweedie objective (variance power 1.1, the v1 reference; and 1.5),
squared error, absolute error (two variants: 31 leaves; 15 leaves with
min. child 100), quantile objective at 0.5, XGBoost {[}Chen \& Guestrin,
2016{]} with absolute error, and ridge regression. Target functionals:
squared error and Tweedie (log link) estimate conditional means;
absolute error and the 0.5 quantile estimate the conditional median.
CatBoost was excluded for cost; statistical models beyond the baselines
(ETS/ARIMA/Theta) are not evaluated because the corresponding packages
were not installed and installing them was not authorised.

\emph{Selection.} For each endpoint the candidate with the lowest mean
validation WAPE of that endpoint (item-day or subset-day) over the three
validation blocks is selected, ties by name; selection is written to
disk before any test block is fitted. Test fits are made only for the
two selected models and the reference model (seeds 43 and 44
additionally for the item-day model on W2).

\emph{Baselines (31).} Sixteen estimators --- naive, seasonal naive,
weekday mean and median (4 and 8 weeks), trailing mean (28, 91 days),
trailing median (14, 28, 56, 91), Croston {[}Croston, 1972{]},
Croston-SBA {[}Syntetos \& Boylan, 2005{]}, TSB {[}Teunter et al.,
2011{]} and an all-zero forecast --- each with \textbf{matched}
information (history up to t − H, exactly the models' information) and
\textbf{fresh} information (re-estimated at the start of each block or
week; 1 to H days old). Baselines are not tuned beyond the choice among
these prespecified variants; the best is selected on validation data for
the same endpoint.

\subsection{7. Endpoints and statistical
analysis}\label{endpoints-and-statistical-analysis}

\emph{Endpoints.} Primary: pooled \textbf{item-day WAPE} = Σ\textbar ŷ −
y\textbar{} / Σy over eligible (series, day) rows. Aggregate:
\textbf{subset-day WAPE} = Σ\_d \textbar Σ\_s (ŷ − y)\textbar{} / Σ\_d
Σ\_s y (demand of the selected subset summed per day before scoring; not
a store total). Descriptive: item-week and subset-week WAPE, relative
bias, MAE, RMSE, share of zero forecasts. Subset-day WAPE is at least
\textbar relative bias\textbar. The estimand is the relative improvement
(E\_ref − E\_model)/E\_ref pooled over the six test blocks; the minimum
relevant improvement is 2 \%.

\emph{Hypotheses.} H1 (primary): the item-day-selected model beats the
best matched baseline on item-day WAPE. H1b: the same against the best
baseline of any information set. H2: the same model beats the v1
reference model (Tweedie 1.1, refit on the same window). H3: the
subset-day-selected model beats the best matched baseline on subset-day
WAPE. H4 (descriptive): endpoint dependence.

\emph{Inference.} Item-day endpoints: paired crossed bootstrap over
series and 7-day blocks (B = 10,000; primary) and a two-way
cluster-robust delta-method interval with a t(G − 1) reference
{[}Cameron et al., 2011{]} (secondary); the crossed bootstrap is the
pigeonhole bootstrap of Owen {[}2007{]} and Owen \& Eckles {[}2012{]}.
Aggregate endpoint: bootstrap over 7-day blocks of daily totals (iid
weeks primary; stationary block bootstrap {[}Politis \& Romano, 1994{]}
secondary; Diebold--Mariano-type tests {[}Diebold \& Mariano, 1995{]}
are an alternative that was not used). The primary methods were fixed by
a pre-freeze simulation (worst-case coverage error 0.050 for the crossed
bootstrap, 0.057 for the iid-week bootstrap; both mildly liberal; the v1
series-only design collapses). Holm correction within each family across
the eight W2 cells. \textbf{Decision per cell:} supported iff Holm p
\textless{} 0.05, the interval excludes 0, improvement ≥ 2 \% and the
model wins ≥ 4 of 6 blocks; contradicted iff the interval lies below 0;
otherwise inconclusive. \textbf{Dual-method rule:} a verdict stands only
if the secondary method gives the same verdict; otherwise it is
inconclusive. Replication sets are analysed descriptively with their own
Holm adjustment and are not pooled with the confirmatory family.

\subsection{8. Results}\label{results}

Every table in this section is generated from the stored analysis files
by \texttt{tools/build\_tables.py} (source files are named under each
table); figures are generated by \texttt{report.py}. Improvements are
relative reductions in WAPE pooled over the six test blocks, (E\_ref −
E\_model)/E\_ref. Intervals are 95 \% percentile intervals from the
crossed bootstrap (item-day) or the iid-week bootstrap (subset-day).
``Holm p'' is adjusted within the eight confirmatory cells of the
window. \textbf{Labels used throughout:} \emph{confirmatory} = window W2
with the frozen rules; \emph{replication} = W3, and W1 for the fresh
store WI\_1; \emph{non-independent re-evaluation} = W1 for CA\_1, TX\_1
and Online Retail II; \emph{descriptive} = prespecified summaries
without a test; \emph{exploratory} = post-hoc or uncontrolled for
multiplicity; \emph{supplementary} = post-hoc calculations added after
the results were known, which do not replace any frozen result.

\subsubsection{8.1 Design as executed}\label{design-as-executed}

\textbf{Table 1. Windows and samples} (validation + test spans from
\texttt{DATA\_AUDIT.md}; rows = series × 168 test days).

{\def\LTcaptype{none} % do not increment counter
\begin{longtable}[]{@{}
  >{\raggedright\arraybackslash}p{(\linewidth - 10\tabcolsep) * \real{0.1667}}
  >{\raggedright\arraybackslash}p{(\linewidth - 10\tabcolsep) * \real{0.1667}}
  >{\raggedright\arraybackslash}p{(\linewidth - 10\tabcolsep) * \real{0.1667}}
  >{\raggedright\arraybackslash}p{(\linewidth - 10\tabcolsep) * \real{0.1667}}
  >{\raggedright\arraybackslash}p{(\linewidth - 10\tabcolsep) * \real{0.1667}}
  >{\raggedright\arraybackslash}p{(\linewidth - 10\tabcolsep) * \real{0.1667}}@{}}
\toprule\noalign{}
\begin{minipage}[b]{\linewidth}\raggedright
window
\end{minipage} & \begin{minipage}[b]{\linewidth}\raggedright
role
\end{minipage} & \begin{minipage}[b]{\linewidth}\raggedright
datasets
\end{minipage} & \begin{minipage}[b]{\linewidth}\raggedright
validation + test span
\end{minipage} & \begin{minipage}[b]{\linewidth}\raggedright
series per dataset
\end{minipage} & \begin{minipage}[b]{\linewidth}\raggedright
rows per cell
\end{minipage} \\
\midrule\noalign{}
\endhead
\bottomrule\noalign{}
\endlastfoot
W2 & \textbf{confirmatory} & M5 CA\_1, TX\_1, WI\_1; Online Retail II &
M5: 2015-01-05 → 2015-09-13; Online Retail II: 2010-07-24 → 2011-04-01 &
350 / 350 / 349 / 100 & 58,800 / 58,800 / 58,632 / 16,800 \\
W3 & replication (independent window) & M5 CA\_1, TX\_1, WI\_1 (not
feasible for Online Retail II) & 2014-04-28 → 2015-01-04 & 350 / 350 /
349 & as W2 \\
W1 & replication for WI\_1 (fresh store); non-independent re-evaluation
for CA\_1, TX\_1, Online Retail II & all four & M5: 2015-09-14 →
2016-05-22; Online Retail II: 2011-04-02 → 2011-12-09 & as W2 & as W2 \\
\end{longtable}
}

Each cell is a dataset × horizon (H = 28 or 7); W2 and W1 have eight
cells, W3 six. Zero shares among test rows were 0.53--0.65 for M5 and
0.32--0.40 for Online Retail II (\texttt{DATA\_AUDIT.md}).

\subsubsection{8.2 Model selection on validation
data}\label{model-selection-on-validation-data}

For the item-day endpoint the absolute-error LightGBM was selected in
all eight W2 cells (31-leaf configuration on M5, the 15-leaf /
minimum-child-100 configuration on Online Retail II). For the subset-day
endpoint the Tweedie (p = 1.1) model was selected in six W2 cells and
the squared-error model in the two WI\_1 cells; the absolute-error and
quantile models were never selected for the aggregate endpoint in any
window. In W1 the 0.5-quantile model replaced the absolute-error model
in two cells (CA\_1 H28, TX\_1 H7). The best validation baseline for
item-day was a trailing or weekday median (trailing 14--91 days; weekday
median of 8 weeks on Online Retail II); for subset-day it was a weekday
mean (4 or 8 weeks). No cell selected the all-zero forecast.

\subsubsection{8.3 Confirmatory results (window
W2)}\label{confirmatory-results-window-w2}

\textbf{Table 2. Item-day endpoint, window W2 (confirmatory):
item-day-selected model (absolute-error LightGBM in all cells) versus
the best matched baseline}

{\def\LTcaptype{none} % do not increment counter
\begin{longtable}[]{@{}
  >{\raggedright\arraybackslash}p{(\linewidth - 16\tabcolsep) * \real{0.1111}}
  >{\raggedright\arraybackslash}p{(\linewidth - 16\tabcolsep) * \real{0.1111}}
  >{\raggedright\arraybackslash}p{(\linewidth - 16\tabcolsep) * \real{0.1111}}
  >{\raggedright\arraybackslash}p{(\linewidth - 16\tabcolsep) * \real{0.1111}}
  >{\raggedright\arraybackslash}p{(\linewidth - 16\tabcolsep) * \real{0.1111}}
  >{\raggedright\arraybackslash}p{(\linewidth - 16\tabcolsep) * \real{0.1111}}
  >{\raggedright\arraybackslash}p{(\linewidth - 16\tabcolsep) * \real{0.1111}}
  >{\raggedright\arraybackslash}p{(\linewidth - 16\tabcolsep) * \real{0.1111}}
  >{\raggedright\arraybackslash}p{(\linewidth - 16\tabcolsep) * \real{0.1111}}@{}}
\toprule\noalign{}
\begin{minipage}[b]{\linewidth}\raggedright
cell
\end{minipage} & \begin{minipage}[b]{\linewidth}\raggedright
best matched baseline
\end{minipage} & \begin{minipage}[b]{\linewidth}\raggedright
WAPE baseline / model
\end{minipage} & \begin{minipage}[b]{\linewidth}\raggedright
improvement {[}95 \% CI{]}
\end{minipage} & \begin{minipage}[b]{\linewidth}\raggedright
Holm p
\end{minipage} & \begin{minipage}[b]{\linewidth}\raggedright
blocks won
\end{minipage} & \begin{minipage}[b]{\linewidth}\raggedright
H1
\end{minipage} & \begin{minipage}[b]{\linewidth}\raggedright
H1b: vs best any-information baseline
\end{minipage} & \begin{minipage}[b]{\linewidth}\raggedright
H2: vs original Tweedie 1.1 model
\end{minipage} \\
\midrule\noalign{}
\endhead
\bottomrule\noalign{}
\endlastfoot
CA\_1 H28 & trailing median 56 & 0.775 / 0.739 & +4.6 \% {[}+2.3,
+7.1{]} & 0.0016 & 6/6 & supported & inconclusive (−1.4 \%) & supported
(+8.2 \%) \\
CA\_1 H7 & trailing median 14 & 0.714 / 0.686 & +4.0 \% {[}+2.3, +6.0{]}
& 0.0016 & 6/6 & supported & supported (+2.5 \%) & supported (+8.0
\%) \\
TX\_1 H28 & trailing median 56 & 0.774 / 0.755 & +2.5 \% {[}+1.0,
+4.1{]} & 0.0036 & 6/6 & supported & inconclusive (+0.8 \%) & supported
(+10.5 \%) \\
TX\_1 H7 & trailing median 91 & 0.752 / 0.718 & +4.5 \% {[}+2.5, +6.8{]}
& 0.0016 & 6/6 & supported & supported (+3.8 \%) & supported (+10.8
\%) \\
WI\_1 H28 & trailing median 56 & 0.783 / 0.757 & +3.3 \% {[}+1.6,
+5.3{]} & 0.0018 & 6/6 & supported & inconclusive (+0.9 \%) & supported
(+7.9 \%) \\
WI\_1 H7 & trailing median 91 & 0.753 / 0.708 & +6.0 \% {[}+3.4, +9.0{]}
& 0.0016 & 6/6 & supported & supported (+3.2 \%) & supported (+8.2
\%) \\
Online Retail II H28 & weekday median 8 & 0.931 / 0.877 & +5.8 \%
{[}+2.7, +9.5{]} & 0.0016 & 6/6 & supported & inconclusive (+2.2 \%) &
supported (+16.7 \%) \\
Online Retail II H7 & weekday median 8 & 0.871 / 0.836 & +4.0 \%
{[}+1.5, +7.0{]} & 0.0036 & 5/6 & supported & supported (+4.0 \%) &
supported (+16.8 \%) \\
\end{longtable}
}

\emph{Source: \texttt{results/analysis/analysis\_W2.json} (H1, H1b, H2);
CSV: \texttt{tables/table2\_h1\_W2.csv}.}

H1 (item-day, model vs best matched baseline) is supported in 8 of 8
cells under the frozen rule (Holm p \textless{} 0.05, interval excluding
zero, improvement ≥ 2 \%, ≥ 4 of 6 blocks won, and the same verdict from
the secondary method). The median improvement is +4.3 \% (range +2.5 \%
\ldots{} +6.0 \%); the model won all six blocks in seven cells and five
of six in Online Retail II H7. The lowest Holm p-values equal the
resolution of the bootstrap (B = 10,000) and are bounds, not exact
values. The smallest effect (TX\_1 H28, +2.5 \%) is close to the 2 \%
threshold. H2 (objective effect): the item-day-selected model improves
on the original Tweedie 1.1 model in 8 of 8 cells (+7.9 \% \ldots{}
+16.8 \%). H1b --- against baselines that are re-estimated at block
starts and therefore see up to 28 days more data than the model --- is
supported in the four H = 7 cells and inconclusive in the four H = 28
cells (point estimates −1.4 \% \ldots{} +2.2 \%); \textbf{the statement
``the learned model beats the best simple baseline'' is therefore
established for baselines with the same information as the model, not
for baselines with fresher information at H = 28.}

\textbf{Table 3. Subset-day endpoint, window W2 (confirmatory):
subset-day-selected model versus the best matched baseline}

{\def\LTcaptype{none} % do not increment counter
\begin{longtable}[]{@{}
  >{\raggedright\arraybackslash}p{(\linewidth - 16\tabcolsep) * \real{0.1111}}
  >{\raggedright\arraybackslash}p{(\linewidth - 16\tabcolsep) * \real{0.1111}}
  >{\raggedright\arraybackslash}p{(\linewidth - 16\tabcolsep) * \real{0.1111}}
  >{\raggedright\arraybackslash}p{(\linewidth - 16\tabcolsep) * \real{0.1111}}
  >{\raggedright\arraybackslash}p{(\linewidth - 16\tabcolsep) * \real{0.1111}}
  >{\raggedright\arraybackslash}p{(\linewidth - 16\tabcolsep) * \real{0.1111}}
  >{\raggedright\arraybackslash}p{(\linewidth - 16\tabcolsep) * \real{0.1111}}
  >{\raggedright\arraybackslash}p{(\linewidth - 16\tabcolsep) * \real{0.1111}}
  >{\raggedright\arraybackslash}p{(\linewidth - 16\tabcolsep) * \real{0.1111}}@{}}
\toprule\noalign{}
\begin{minipage}[b]{\linewidth}\raggedright
cell
\end{minipage} & \begin{minipage}[b]{\linewidth}\raggedright
model
\end{minipage} & \begin{minipage}[b]{\linewidth}\raggedright
best matched baseline
\end{minipage} & \begin{minipage}[b]{\linewidth}\raggedright
WAPE baseline / model
\end{minipage} & \begin{minipage}[b]{\linewidth}\raggedright
improvement {[}95 \% CI{]}
\end{minipage} & \begin{minipage}[b]{\linewidth}\raggedright
frozen Holm p primary / secondary
\end{minipage} & \begin{minipage}[b]{\linewidth}\raggedright
blocks won
\end{minipage} & \begin{minipage}[b]{\linewidth}\raggedright
verdict (frozen rule)
\end{minipage} & \begin{minipage}[b]{\linewidth}\raggedright
supplementary converged Holm p primary / secondary
\end{minipage} \\
\midrule\noalign{}
\endhead
\bottomrule\noalign{}
\endlastfoot
CA\_1 H28 & Tweedie 1.1 & weekday mean 4 & 0.114 / 0.088 & +22.7 \%
{[}+14.6, +30.6{]} & 0.002 / 0.002 & 6/6 & supported & \textless0.001 /
\textless0.001 \\
CA\_1 H7 & Tweedie 1.1 & weekday mean 8 & 0.104 / 0.084 & +18.7 \%
{[}+9.5, +28.5{]} & 0.002 / 0.002 & 6/6 & supported & \textless0.001 /
\textless0.001 \\
TX\_1 H28 & Tweedie 1.1 & weekday mean 8 & 0.117 / 0.099 & +14.8 \%
{[}+2.5, +26.7{]} & 0.071 / 0.048 & 4/6 & inconclusive & 0.051 /
0.043 \\
TX\_1 H7 & Tweedie 1.1 & weekday mean 4 & 0.123 / 0.094 & +23.8 \%
{[}+13.9, +33.1{]} & 0.002 / 0.002 & 6/6 & supported & \textless0.001 /
\textless0.001 \\
WI\_1 H28 & squared error & weekday mean 8 & 0.100 / 0.092 & +7.8 \%
{[}−0.5, +15.7{]} & 0.205 / 0.269 & 5/6 & inconclusive & 0.215 /
0.259 \\
WI\_1 H7 & squared error & weekday mean 8 & 0.099 / 0.087 & +11.4 \%
{[}+3.6, +18.3{]} & 0.035 / 0.002 & 6/6 & supported & 0.029 / 0.001 \\
Online Retail II H28 & Tweedie 1.1 & weekday mean 8 & 0.511 / 0.482 &
+5.8 \% {[}−5.5, +16.2{]} & 0.582 / 0.405 & 4/6 & inconclusive & 0.596 /
0.403 \\
Online Retail II H7 & Tweedie 1.1 & weekday mean 8 & 0.456 / 0.458 &
−0.5 \% {[}−13.3, +8.4{]} & 0.939 / 0.934 & 2/6 & inconclusive & 0.938 /
0.927 \\
\end{longtable}
}

\emph{Source: \texttt{results/analysis/analysis\_W2.json} (H3);
\texttt{results/supplementary/h3\_convergence.json}; CSV:
\texttt{tables/table3\_h3\_W2.csv}.}

H3 (subset-day, model vs best matched baseline) is supported in 4 of 8
cells and inconclusive in 4; no cell is contradicted. Point estimates
are positive in 7 of 8 cells (median +13.1 \%; Online Retail II H7 −0.5
\%). The intervals are wide because the aggregate endpoint has only 24
weekly blocks. TX\_1 H28 is a boundary case examined in Section 8.5; its
frozen verdict (inconclusive) is confirmed.

\subsubsection{8.4 Replication windows}\label{replication-windows}

\textbf{Table 4. Verdict counts (supported / contradicted /
inconclusive) by window; each window has its own Holm adjustment}

{\def\LTcaptype{none} % do not increment counter
\begin{longtable}[]{@{}
  >{\raggedright\arraybackslash}p{(\linewidth - 14\tabcolsep) * \real{0.1250}}
  >{\raggedright\arraybackslash}p{(\linewidth - 14\tabcolsep) * \real{0.1250}}
  >{\raggedright\arraybackslash}p{(\linewidth - 14\tabcolsep) * \real{0.1250}}
  >{\raggedright\arraybackslash}p{(\linewidth - 14\tabcolsep) * \real{0.1250}}
  >{\raggedright\arraybackslash}p{(\linewidth - 14\tabcolsep) * \real{0.1250}}
  >{\raggedright\arraybackslash}p{(\linewidth - 14\tabcolsep) * \real{0.1250}}
  >{\raggedright\arraybackslash}p{(\linewidth - 14\tabcolsep) * \real{0.1250}}
  >{\raggedright\arraybackslash}p{(\linewidth - 14\tabcolsep) * \real{0.1250}}@{}}
\toprule\noalign{}
\begin{minipage}[b]{\linewidth}\raggedright
window
\end{minipage} & \begin{minipage}[b]{\linewidth}\raggedright
cells
\end{minipage} & \begin{minipage}[b]{\linewidth}\raggedright
H1 s/c/i
\end{minipage} & \begin{minipage}[b]{\linewidth}\raggedright
H1b s/c/i
\end{minipage} & \begin{minipage}[b]{\linewidth}\raggedright
H2 s/c/i
\end{minipage} & \begin{minipage}[b]{\linewidth}\raggedright
H3 s/c/i
\end{minipage} & \begin{minipage}[b]{\linewidth}\raggedright
median H1 improvement (range)
\end{minipage} & \begin{minipage}[b]{\linewidth}\raggedright
H1 improvement \textgreater{} 0
\end{minipage} \\
\midrule\noalign{}
\endhead
\bottomrule\noalign{}
\endlastfoot
W2 (confirmatory) & 8 & 8/0/0 & 4/0/4 & 8/0/0 & 4/0/4 & +4.3 \% (+2.5 \%
\ldots{} +6.0 \%) & 8/8 \\
W3 (M5, independent replication) & 6 & 6/0/0 & 3/0/3 & 6/0/0 & 4/0/2 &
+4.5 \% (+2.2 \% \ldots{} +6.2 \%) & 6/6 \\
W1 (WI\_1 fresh store; CA\_1, TX\_1, Online Retail II non-independent) &
8 & 6/0/2 & 3/0/5 & 8/0/0 & 4/0/4 & +3.2 \% (+1.1 \% \ldots{} +4.5 \%) &
8/8 \\
\end{longtable}
}

\emph{Source: \texttt{results/summary\_extra.json} (from
\texttt{analysis\_*.json}); CSV:
\texttt{tables/table4\_replication.csv}.}

The H1 improvement is positive in all 22 cells of the three windows
(supported in 20). In W3 it is supported in all six cells. In W1 two
cells are inconclusive: TX\_1 H28 (+1.1 \%, 3/6 blocks won, interval
includes zero) and TX\_1 H7 (+1.9 \%, below the 2 \% threshold although
the interval excludes zero); the two WI\_1 cells of W1 (the fresh store)
are supported. H3 is supported in 4 of 6 cells in W3 and in 4 of 8 in
W1. Per-cell results are in Table A1 (Appendix) and Figure 1. W1 results
for CA\_1, TX\_1 and Online Retail II are not independent evidence
(Section 5).

\begin{figure}
\centering
\pandocbounded{\includegraphics[keepaspectratio,alt={Figure 1. Relative improvement in WAPE with 95 \% intervals (crossed bootstrap for H1 and H2, iid-week bootstrap for H3), by window (blue W2 confirmatory, orange W3 replication, grey W1; only WI\_1 is a fresh replication in W1). Dotted line: 2 \% minimum improvement. Source: results/analysis/analysis\_*.json, generated by report.py.}]{figures/fig1_forest.png}}
\caption{Figure 1. Relative improvement in WAPE with 95 \% intervals
(crossed bootstrap for H1 and H2, iid-week bootstrap for H3), by window
(blue W2 confirmatory, orange W3 replication, grey W1; only WI\_1 is a
fresh replication in W1). Dotted line: 2 \% minimum improvement. Source:
\protect\texttt{results/analysis/analysis\_*.json}, generated by
\protect\texttt{report.py}.}
\end{figure}

\emph{Figure 1. Relative improvement in WAPE with 95 \% intervals
(crossed bootstrap for H1 and H2, iid-week bootstrap for H3), by window
(blue W2 confirmatory, orange W3 replication, grey W1; only WI\_1 is a
fresh replication in W1). Dotted line: 2 \% minimum improvement. Source:
\texttt{results/analysis/analysis\_*.json}, generated by
\texttt{report.py}.}

\subsubsection{8.5 Supplementary precision check of the boundary cell
(TX\_1, H = 28,
subset-day)}\label{supplementary-precision-check-of-the-boundary-cell-tx_1-h-28-subset-day}

The independent re-implementation (B = 4,000) classified this cell as
supported whereas the frozen analysis (B = 10,000, fixed seed)
classified it as inconclusive because the two prespecified methods
disagreed. \textbf{This section is supplementary: it was added after the
results were known, uses the same two prespecified methods and the
unchanged decision rule, and does not replace the frozen result.}

\textbf{Table 6. W2, TX\_1, H = 28, subset-day (H3): frozen result and
supplementary convergence check}

{\def\LTcaptype{none} % do not increment counter
\begin{longtable}[]{@{}
  >{\raggedright\arraybackslash}p{(\linewidth - 4\tabcolsep) * \real{0.3333}}
  >{\raggedright\arraybackslash}p{(\linewidth - 4\tabcolsep) * \real{0.3333}}
  >{\raggedright\arraybackslash}p{(\linewidth - 4\tabcolsep) * \real{0.3333}}@{}}
\toprule\noalign{}
\begin{minipage}[b]{\linewidth}\raggedright
quantity
\end{minipage} & \begin{minipage}[b]{\linewidth}\raggedright
frozen run (B = 10,000)
\end{minipage} & \begin{minipage}[b]{\linewidth}\raggedright
supplementary (3 × 10\^{}6 draws)
\end{minipage} \\
\midrule\noalign{}
\endhead
\bottomrule\noalign{}
\endlastfoot
raw p, primary (iid-week bootstrap) & 0.0178 & 0.0129 (seed SD
0.0002) \\
raw p, secondary (stationary bootstrap) & 0.0120 & 0.0108 (seed SD
0.0000) \\
Holm p, primary & 0.0712 & 0.0514 \\
Holm p, secondary & 0.0480 & 0.0433 \\
verdict under the frozen dual-method rule & inconclusive &
inconclusive \\
exact sign-flip test (unadjusted p; different statistic) & -- &
0.0339 \\
weeks in which the model has the smaller absolute total error & -- &
15/24 \\
lag-1 autocorrelation of the weekly differences / AR(1) effective weeks
& -- & +0.36 / 11.3 \\
\end{longtable}
}

\emph{Source: \texttt{results/analysis/analysis\_W2.json};
\texttt{results/supplementary/h3\_convergence.json}; CSV:
\texttt{tables/table6\_tx1\_h28\_h3.csv}.}

With 3 × 10\^{}6 draws per method (three new seeds) Monte-Carlo error is
negligible (seed SD ≤ 0.0002). The primary method's Holm p is 0.0514 and
the secondary method's 0.0433: the methods still disagree and the frozen
rule gives \emph{inconclusive}. The frozen verdict is confirmed, and the
independent code's \emph{supported} was Monte-Carlo noise at B = 4,000.
All 22 H3 verdicts of the three windows are unchanged under the
converged calculation (\texttt{SUPPLEMENTARY\_H3\_RESOLUTION.md}). The
weekly differences in this cell are positively autocorrelated (lag-1
+0.36; AR(1) effective sample ≈ 11 of 24 weeks), which the iid-week
bootstrap ignores; the cell is best described as suggestive, not
established.

\subsubsection{8.6 Seed sensitivity (W2,
descriptive)}\label{seed-sensitivity-w2-descriptive}

Refitting the item-day model with seeds 43 and 44 changes the H1
improvement by at most 0.4 percentage points (per cell 0.09--0.40 pp;
Table A2); the lowest of the three seeds is ≥ 2.3 \% in every cell. Seed
variation is not included in the intervals.

\subsubsection{8.7 Aggregation and bias (descriptive, prespecified
H4)}\label{aggregation-and-bias-descriptive-prespecified-h4}

\textbf{Table 5. Aggregation ladder and bias, window W2, all cells}

{\def\LTcaptype{none} % do not increment counter
\begin{longtable}[]{@{}
  >{\raggedright\arraybackslash}p{(\linewidth - 14\tabcolsep) * \real{0.1250}}
  >{\raggedright\arraybackslash}p{(\linewidth - 14\tabcolsep) * \real{0.1250}}
  >{\raggedright\arraybackslash}p{(\linewidth - 14\tabcolsep) * \real{0.1250}}
  >{\raggedright\arraybackslash}p{(\linewidth - 14\tabcolsep) * \real{0.1250}}
  >{\raggedright\arraybackslash}p{(\linewidth - 14\tabcolsep) * \real{0.1250}}
  >{\raggedright\arraybackslash}p{(\linewidth - 14\tabcolsep) * \real{0.1250}}
  >{\raggedright\arraybackslash}p{(\linewidth - 14\tabcolsep) * \real{0.1250}}
  >{\raggedright\arraybackslash}p{(\linewidth - 14\tabcolsep) * \real{0.1250}}@{}}
\toprule\noalign{}
\begin{minipage}[b]{\linewidth}\raggedright
cell
\end{minipage} & \begin{minipage}[b]{\linewidth}\raggedright
method
\end{minipage} & \begin{minipage}[b]{\linewidth}\raggedright
item-day WAPE
\end{minipage} & \begin{minipage}[b]{\linewidth}\raggedright
item-week WAPE
\end{minipage} & \begin{minipage}[b]{\linewidth}\raggedright
subset-week WAPE
\end{minipage} & \begin{minipage}[b]{\linewidth}\raggedright
subset-day WAPE
\end{minipage} & \begin{minipage}[b]{\linewidth}\raggedright
relative bias
\end{minipage} & \begin{minipage}[b]{\linewidth}\raggedright
share of zero forecasts
\end{minipage} \\
\midrule\noalign{}
\endhead
\bottomrule\noalign{}
\endlastfoot
CA\_1 H28 & absolute-error model (item-day winner) & 0.739 & 0.518 &
0.346 & 0.346 & −0.346 & 0.29 \\
CA\_1 H28 & best item-day baseline & 0.775 & 0.560 & 0.293 & 0.295 &
−0.293 & 0.57 \\
CA\_1 H28 & Tweedie 1.1 model (subset-day winner) & 0.806 & 0.441 &
0.046 & 0.088 & +0.001 & 0.00 \\
CA\_1 H28 & best subset-day baseline & 0.849 & 0.454 & 0.074 & 0.114 &
−0.007 & 0.24 \\
CA\_1 H7 & absolute-error model (item-day winner) & 0.686 & 0.429 &
0.287 & 0.287 & −0.287 & 0.29 \\
CA\_1 H7 & best item-day baseline & 0.714 & 0.448 & 0.232 & 0.242 &
−0.232 & 0.53 \\
CA\_1 H7 & Tweedie 1.1 model (subset-day winner) & 0.745 & 0.368 & 0.045
& 0.084 & −0.004 & 0.00 \\
CA\_1 H7 & best subset-day baseline & 0.792 & 0.405 & 0.068 & 0.104 &
−0.006 & 0.13 \\
TX\_1 H28 & absolute-error model (item-day winner) & 0.755 & 0.558 &
0.394 & 0.394 & −0.394 & 0.49 \\
TX\_1 H28 & best item-day baseline & 0.774 & 0.573 & 0.330 & 0.332 &
−0.330 & 0.68 \\
TX\_1 H28 & Tweedie 1.1 model (subset-day winner) & 0.843 & 0.464 &
0.060 & 0.099 & −0.021 & 0.00 \\
TX\_1 H28 & best subset-day baseline & 0.881 & 0.490 & 0.074 & 0.117 &
+0.010 & 0.19 \\
TX\_1 H7 & absolute-error model (item-day winner) & 0.718 & 0.487 &
0.360 & 0.360 & −0.360 & 0.48 \\
TX\_1 H7 & best item-day baseline & 0.752 & 0.552 & 0.346 & 0.347 &
−0.346 & 0.70 \\
TX\_1 H7 & Tweedie 1.1 model (subset-day winner) & 0.805 & 0.406 & 0.054
& 0.094 & −0.019 & 0.00 \\
TX\_1 H7 & best subset-day baseline & 0.862 & 0.428 & 0.071 & 0.123 &
+0.005 & 0.32 \\
WI\_1 H28 & absolute-error model (item-day winner) & 0.757 & 0.505 &
0.331 & 0.331 & −0.331 & 0.35 \\
WI\_1 H28 & best item-day baseline & 0.783 & 0.534 & 0.289 & 0.290 &
−0.289 & 0.59 \\
WI\_1 H28 & squared error model (subset-day winner) & 0.831 & 0.433 &
0.047 & 0.092 & +0.020 & 0.00 \\
WI\_1 H28 & best subset-day baseline & 0.855 & 0.450 & 0.052 & 0.100 &
+0.003 & 0.17 \\
WI\_1 H7 & absolute-error model (item-day winner) & 0.708 & 0.427 &
0.307 & 0.307 & −0.307 & 0.38 \\
WI\_1 H7 & best item-day baseline & 0.753 & 0.502 & 0.298 & 0.298 &
−0.298 & 0.60 \\
WI\_1 H7 & squared error model (subset-day winner) & 0.784 & 0.370 &
0.043 & 0.087 & +0.005 & 0.00 \\
WI\_1 H7 & best subset-day baseline & 0.822 & 0.400 & 0.045 & 0.099 &
+0.000 & 0.17 \\
Online Retail II H28 & absolute-error model (item-day winner) & 0.877 &
0.662 & 0.515 & 0.573 & −0.384 & 0.07 \\
Online Retail II H28 & best item-day baseline & 0.931 & 0.691 & 0.500 &
0.572 & −0.329 & 0.29 \\
Online Retail II H28 & Tweedie 1.1 model (subset-day winner) & 1.053 &
0.706 & 0.312 & 0.482 & +0.085 & 0.00 \\
Online Retail II H28 & best subset-day baseline & 1.135 & 0.765 & 0.402
& 0.511 & +0.115 & 0.20 \\
Online Retail II H7 & absolute-error model (item-day winner) & 0.836 &
0.606 & 0.490 & 0.534 & −0.390 & 0.03 \\
Online Retail II H7 & best item-day baseline & 0.871 & 0.618 & 0.483 &
0.541 & −0.361 & 0.29 \\
Online Retail II H7 & Tweedie 1.1 model (subset-day winner) & 1.005 &
0.649 & 0.305 & 0.458 & +0.059 & 0.00 \\
Online Retail II H7 & best subset-day baseline & 1.057 & 0.665 & 0.292 &
0.456 & +0.068 & 0.20 \\
\end{longtable}
}

\emph{Source: \texttt{results/analysis/analysis\_W2.json}
(\texttt{ladder}); CSV: \texttt{tables/table5\_ladder\_W2.csv}.}

The model that wins the item-day endpoint is strongly downward-biased
(relative bias −0.29 \ldots{} −0.39 on M5; it forecasts zero for 29--49
\% of M5 rows against 53--70 \% for the best item-day baseline). Because
subset-day WAPE is at least the absolute relative bias (with equality
for this model on M5), the same model loses to simple baselines on the
aggregate endpoint (for example CA\_1 H28: 0.346 against 0.114 for the
best subset-day baseline), and the models selected for the two endpoints
differ in all eight cells. The original reference model (Tweedie 1.1)
has bias near zero and loses to the best baseline on the item-day
endpoint (CA\_1 H28: 0.806 against 0.775).

\begin{figure}
\centering
\pandocbounded{\includegraphics[keepaspectratio,alt={Figure 2. Aggregation ladder, window W2: WAPE of the item-day-selected model and of the best matched item-day baseline at four aggregation levels. Source: results/analysis/analysis\_W2.json (ladder); Table 5 gives all values including the subset-day-selected models.}]{figures/fig2_aggregation_ladder_W2.png}}
\caption{Figure 2. Aggregation ladder, window W2: WAPE of the
item-day-selected model and of the best matched item-day baseline at
four aggregation levels. Source:
\protect\texttt{results/analysis/analysis\_W2.json}
(\protect\texttt{ladder}); Table 5 gives all values including the
subset-day-selected models.}
\end{figure}

\emph{Figure 2. Aggregation ladder, window W2: WAPE of the
item-day-selected model and of the best matched item-day baseline at
four aggregation levels. Source:
\texttt{results/analysis/analysis\_W2.json} (\texttt{ladder}); Table 5
gives all values including the subset-day-selected models.}

\begin{figure}
\centering
\pandocbounded{\includegraphics[keepaspectratio,alt={Figure 3. Item-day improvement (H1) against subset-day improvement (H3) per cell and window. These are different models: the item-day and subset-day winners are selected separately. Source: results/analysis/analysis\_*.json.}]{figures/fig3_endpoint_dependence.png}}
\caption{Figure 3. Item-day improvement (H1) against subset-day
improvement (H3) per cell and window. These are different models: the
item-day and subset-day winners are selected separately. Source:
\protect\texttt{results/analysis/analysis\_*.json}.}
\end{figure}

\emph{Figure 3. Item-day improvement (H1) against subset-day improvement
(H3) per cell and window. These are different models: the item-day and
subset-day winners are selected separately. Source:
\texttt{results/analysis/analysis\_*.json}.}

\subsubsection{8.8 Exploratory results}\label{exploratory-results}

Gains of the item-day model are concentrated in frequently selling and
high-volume series (Table A3, W2; for example CA\_1 H28: frequent series
+8.1 \% {[}+3.5, +12.7{]}, n = 46; intermittent or sparse +2.5 \%
{[}+0.5, +4.7{]}, n = 304; low and mid volume terciles with intervals
that include zero). These segment intervals (B = 2,000) are not
corrected for multiplicity and are descriptive. The design effect of the
crossed series × week structure on the standard error is 1.5--2.5 (Table
A4).

\subsection{9. Discussion}\label{discussion}

\textbf{What the study adds.} The negative v1 result was largely an
objective--metric mismatch plus baseline-level bias, not evidence that
boosted models cannot beat simple baselines: when an absolute-error
objective is in the candidate set and selection is made on the same
endpoint, the learned model beats the best same-information baseline on
the item-day endpoint by a modest 2.5--6.0 \% in every confirmatory cell
and in the replication windows, while Tweedie / squared-error models win
the aggregate endpoint by larger margins that are less certain
(supported in 4 of 8 W2 cells). Both findings follow from
scoring-function theory {[}Gneiting, 2011; Kolassa, 2020{]}; the
contribution is their quantification under a frozen protocol with
dependence-aware intervals on public retail data, and the demonstration
that the answer to ``does the learned model beat a simple baseline?''
can flip with the endpoint.

\textbf{What the study does not show.} (i) That the item-day gain is
useful: it comes from a model that forecasts zero more often and has
large negative bias; whether that is desirable depends on a decision
loss that was not evaluated. (ii) That learned models beat baselines
with fresher information at H = 28 (H1b). (iii) That any model is best:
the candidate set is small (eight configurations from LightGBM, XGBoost
and ridge regression), tuning is limited, and statistical models (ETS,
ARIMA, Theta), neural networks and probabilistic forecasts were not
evaluated; Section 10 and \texttt{PROPOSAL\_STATISTICAL\_BASELINES.md}
explain why this matters and how it could be addressed on a fresh
window. (iv) Generality beyond three stores of one retailer and one UK
wholesaler. (v) Causal explanations: the interpretation in terms of bias
and zero forecasts rests on a descriptive decomposition, not an
intervention.

\textbf{Relationship to the literature.} Tree-based models were
prevalent among top M5 methods and outperformed other approaches there
{[}Makridakis et al., 2022; Januschowski et al., 2022{]}; the M5 report
also finds the top-performing methods clearly better than its
statistical benchmarks on its own (hierarchical) metric. The
inconsistency of conventional measures on intermittent demand is
documented {[}Prestwich et al., 2014; Hyndman \& Koehler, 2006{]}. The
present results are consistent with both and add effect sizes for a
matched-information comparison against a large baseline set, with no
claim that M5-style evaluation (recursive, hierarchical, with price and
calendar covariates and a weighted scaled metric) is reproduced.

\textbf{Practical reading (cautious).} The endpoint should be chosen
before models are compared and should reflect the decision. If the
decision depends on totals (for example weekly store volume), a
mean-type objective and an aggregate metric are appropriate; if item-day
accuracy in the absolute-error sense is wanted, a median-type objective
helps, with the bias that implies. No cost, service-level or profit
analysis was done.

\subsection{10. Limitations and threats to
validity}\label{limitations-and-threats-to-validity}

\begin{itemize}
\tightlist
\item
  \textbf{Internal freeze, not public pre-registration.} The protocol
  was hashed before any W2/W3/WI\_1 outcome was read, but the hash is
  not externally timestamped and the design was set with knowledge of
  the v1 results and an internal audit of them.
\item
  \textbf{Window exposure.} W2 and W3 targets of the evaluated entities
  were not scored before; their calendar \emph{periods} were used as
  training data and in a retired benchmark on other M5 series
  (documented). W1 for CA\_1, TX\_1 and Online Retail II is
  non-independent. Online Retail II has no W3 and a short W2 history
  (487 days).
\item
  \textbf{Baselines and models.} 31 prespecified simple baselines and
  eight model configurations with shared hyper-parameters. \textbf{ETS,
  ARIMA, Theta, global neural models and probabilistic or quantile-loss
  evaluations are absent}; statistical-model packages were not installed
  and installing them was not authorised. Readers who consider these
  necessary for the claim ``beats the best simple baseline'' should
  treat the present result as limited to the stated baseline set; a
  separate, fresh-window proposal exists but was not run.
\item
  \textbf{Sampling.} 350/349 department-stratified series per M5 store
  (series launched after d\_1100 excluded) and the 100 largest-revenue
  Online Retail II SKUs; results describe these samples. The eligibility
  rule (days with a listed price) is an assumption about what a
  forecaster knows.
\item
  \textbf{Dependence and inference.} Intervals are percentile or t-based
  and mildly liberal in simulation (worst-case coverage ≈ 0.90--0.95 at
  nominal 0.95); the aggregate endpoint has 24 weekly blocks; the
  iid-week bootstrap ignores weekly autocorrelation (visible in TX\_1
  H28); model errors across blocks are correlated through overlapping
  training windows (not modelled); one seed per primary model, seed
  variation reported separately.
\item
  \textbf{Metric choice.} WAPE and its aggregate version are the
  prespecified endpoints; scaled, probabilistic, cost-based or
  service-level criteria were not evaluated. Item-day WAPE rewards
  forecasts near the conditional median on sparse series.
\item
  \textbf{Dual-method rule.} It converts method disagreement into
  \emph{inconclusive}, which is conservative for H3 (one W2 cell, two W3
  cells); the supplementary check shows the verdict counts do not change
  when Monte-Carlo error is removed.
\item
  \textbf{Same-agent development and verification.} The study code, the
  verification code, the supplementary analysis and this manuscript were
  produced by one AI assistant under the author's direction. The
  verification is an internal consistency and re-implementation check,
  not independent external validation (Section 11).
\item
  \textbf{Data rights.} M5 data and any array containing M5 sales are
  not released.
\end{itemize}

\subsection{11. Verification}\label{verification}

Unit tests (17) cover leak-freedom of all 16 baseline estimators under
both information sets, metric identities, objective semantics, window
arithmetic, the Online Retail II SKU list and the inference helpers;
they pass in the project environment and in a fresh environment with
identical pinned versions. A second code path
(\texttt{verify\_independent.py}), which imports nothing from the study,
recomputed targets from the raw files, all 31 baselines row by row,
metrics, improvements, bootstrap and cluster-robust intervals (own
implementation, B = 4,000), Holm adjustments and verdicts, the
validation-only selection and the ordering of selection before test
fits: 98/98 checks passed in W2, 72/72 in W3 and 98/98 in W1; the frozen
files were unchanged and point estimates were identical in 88 of 88
comparisons. Of the 88 verdict comparisons 87 agree; the exception (W2
TX\_1 H28, H3) is resolved in Section 8.5. One complete cell (Online
Retail II, W2, H = 7) was re-run from scratch in a fresh environment
with identical pinned versions and reproduced all 35 stored method
summaries exactly (maximum absolute difference 0); the other 21 cells
were not re-run. The same AI assistant designed and ran all of this;
\textbf{it is not external validation}. A human reviewer should in
particular check the data-licence reading, the reference list against
original sources, the interpretation of the boundary cell, and, ideally,
replicate the pipeline independently.

\subsection{12. Conclusions}\label{conclusions}

In the confirmatory window, an absolute-error LightGBM selected on
validation item-day WAPE beat the best simple baseline with the same
information in all eight dataset--horizon cells (median +4.3 \%, range
+2.5 \ldots{} +6.0 \%); the improvement was positive in all 22 cells of
three windows (supported in 20) and the model beat the earlier Tweedie
reference in every cell. Against baselines with fresher information at H
= 28 the advantage is not established, and statistical, neural and
probabilistic competitors were not evaluated. On the aggregate endpoint
a different, mean-type model is selected; the item-day winner loses to
simple baselines because of its level bias, and the improvement of the
aggregate-selected model over the best baseline is supported in four of
eight confirmatory cells (positive in seven), with the TX\_1 H28 cell a
documented boundary case. The take-home message is endpoint-matching:
the same data and models give opposite answers depending on the metric,
so a claim that a model beats a baseline is meaningful only relative to
the endpoint, the information set, the baseline set and the baseline's
level bias.

\subsection{13. Reproducibility, data and code
availability}\label{reproducibility-data-and-code-availability}

Code, frozen protocol, configuration, result tables and figures are in
\texttt{research/forecast\_study\_v2/} (commands in
\texttt{REPRODUCIBILITY.md}). Raw M5 data must be obtained from Kaggle
under the competition rules (account and acceptance of the rules
required); Online Retail II is available from the UCI repository (CC BY
4.0, doi:10.24432/C5CG6D). Stored per-row predictions contain M5 sales
and are not released; aggregate tables, figures, code, configuration and
hashes can be shared. One cell was reproduced exactly in a fresh
environment; a full clean-room reproduction has not been performed
(\texttt{FINAL\_VERIFICATION\_REPORT.md}).

\subsection{14. Declarations (to be completed by the
authors)}\label{declarations-to-be-completed-by-the-authors}

\begin{itemize}
\tightlist
\item
  \textbf{Author names, affiliations, contributions:} {[}AUTHORS TO
  COMPLETE{]}
\item
  \textbf{Funding:} {[}AUTHORS TO COMPLETE{]}
\item
  \textbf{Competing interests:} {[}AUTHORS TO COMPLETE{]}
\item
  \textbf{Ethics:} the study uses public, secondary, non-interventional
  data; Online Retail II contains pseudonymous customer identifiers but
  no customer-level data were analysed. {[}AUTHORS TO CONFIRM and, if
  required, obtain institutional confirmation{]}
\item
  \textbf{Use of generative AI:} {[}AUTHORS TO COMPLETE.{]} Factual note
  for the authors: the experimental code, analysis and verification
  code, the supplementary analysis and the drafting of this manuscript
  were produced with an AI model (Anthropic Claude, via Claude Code)
  under the author's direction; the authors must review all content,
  take responsibility for it and comply with the venue's AI-use policy.
\item
  \textbf{Data availability:} as in Section 13; M5 data cannot be
  redistributed.
\item
  \textbf{Pre-registration:} none (internal timestamped freeze only).
\end{itemize}

\subsection{Appendix A. Supplementary
tables}\label{appendix-a.-supplementary-tables}

\textbf{Table A1. Replication windows W3 and W1: item-day (H1) and
subset-day (H3) results}

{\def\LTcaptype{none} % do not increment counter
\begin{longtable}[]{@{}
  >{\raggedright\arraybackslash}p{(\linewidth - 18\tabcolsep) * \real{0.1000}}
  >{\raggedright\arraybackslash}p{(\linewidth - 18\tabcolsep) * \real{0.1000}}
  >{\raggedright\arraybackslash}p{(\linewidth - 18\tabcolsep) * \real{0.1000}}
  >{\raggedright\arraybackslash}p{(\linewidth - 18\tabcolsep) * \real{0.1000}}
  >{\raggedright\arraybackslash}p{(\linewidth - 18\tabcolsep) * \real{0.1000}}
  >{\raggedright\arraybackslash}p{(\linewidth - 18\tabcolsep) * \real{0.1000}}
  >{\raggedright\arraybackslash}p{(\linewidth - 18\tabcolsep) * \real{0.1000}}
  >{\raggedright\arraybackslash}p{(\linewidth - 18\tabcolsep) * \real{0.1000}}
  >{\raggedright\arraybackslash}p{(\linewidth - 18\tabcolsep) * \real{0.1000}}
  >{\raggedright\arraybackslash}p{(\linewidth - 18\tabcolsep) * \real{0.1000}}@{}}
\toprule\noalign{}
\begin{minipage}[b]{\linewidth}\raggedright
window
\end{minipage} & \begin{minipage}[b]{\linewidth}\raggedright
cell
\end{minipage} & \begin{minipage}[b]{\linewidth}\raggedright
item-day model
\end{minipage} & \begin{minipage}[b]{\linewidth}\raggedright
best matched baseline
\end{minipage} & \begin{minipage}[b]{\linewidth}\raggedright
WAPE baseline / model
\end{minipage} & \begin{minipage}[b]{\linewidth}\raggedright
H1 improvement {[}95 \% CI{]}
\end{minipage} & \begin{minipage}[b]{\linewidth}\raggedright
Holm p
\end{minipage} & \begin{minipage}[b]{\linewidth}\raggedright
blocks won
\end{minipage} & \begin{minipage}[b]{\linewidth}\raggedright
H1
\end{minipage} & \begin{minipage}[b]{\linewidth}\raggedright
H3 (improvement)
\end{minipage} \\
\midrule\noalign{}
\endhead
\bottomrule\noalign{}
\endlastfoot
W3 & CA\_1 H28 & absolute error & trailing median 91 & 0.771 / 0.732 &
+5.0 \% {[}+2.9, +7.5{]} & 0.0012 & 6/6 & supported & supported (+21.1
\%) \\
W3 & CA\_1 H7 & absolute error & trailing median 14 & 0.734 / 0.689 &
+6.2 \% {[}+4.5, +8.0{]} & 0.0012 & 6/6 & supported & inconclusive
(+16.0 \%) \\
W3 & TX\_1 H28 & absolute error & trailing median 91 & 0.793 / 0.776 &
+2.2 \% {[}+0.2, +4.1{]} & 0.0284 & 5/6 & supported & supported (+18.0
\%) \\
W3 & TX\_1 H7 & absolute error & trailing median 28 & 0.758 / 0.726 &
+4.1 \% {[}+2.4, +6.0{]} & 0.0012 & 6/6 & supported & supported (+19.1
\%) \\
W3 & WI\_1 H28 & absolute error & trailing median 56 & 0.800 / 0.764 &
+4.5 \% {[}+2.7, +6.2{]} & 0.0012 & 6/6 & supported & supported (+13.5
\%) \\
W3 & WI\_1 H7 & absolute error & trailing median 28 & 0.751 / 0.717 &
+4.5 \% {[}+3.1, +6.0{]} & 0.0012 & 6/6 & supported & inconclusive
(+11.7 \%) \\
W1 & CA\_1 H28 & 0.5 quantile & trailing median 28 & 0.724 / 0.708 &
+2.3 \% {[}+0.7, +3.9{]} & 0.0152 & 6/6 & supported & inconclusive
(+12.2 \%) \\
W1 & CA\_1 H7 & absolute error & trailing median 14 & 0.691 / 0.661 &
+4.3 \% {[}+2.9, +5.8{]} & 0.0016 & 6/6 & supported & supported (+18.4
\%) \\
W1 & TX\_1 H28 & absolute error & trailing median 56 & 0.752 / 0.743 &
+1.1 \% {[}−1.0, +3.2{]} & 0.2848 & 3/6 & inconclusive & inconclusive
(+8.3 \%) \\
W1 & TX\_1 H7 & 0.5 quantile & trailing median 28 & 0.710 / 0.697 & +1.9
\% {[}+0.6, +3.3{]} & 0.0138 & 6/6 & inconclusive & inconclusive (+12.7
\%) \\
W1 & WI\_1 H28 & absolute error & trailing median 56 & 0.753 / 0.725 &
+3.8 \% {[}+2.3, +5.4{]} & 0.0016 & 6/6 & supported & supported (+13.5
\%) \\
W1 & WI\_1 H7 & absolute error & trailing median 28 & 0.724 / 0.691 &
+4.5 \% {[}+2.9, +6.2{]} & 0.0016 & 6/6 & supported & supported (+16.4
\%) \\
W1 & Online Retail II H28 & absolute error (small) & weekday median 8 &
0.817 / 0.790 & +3.3 \% {[}+1.7, +5.1{]} & 0.0016 & 6/6 & supported &
supported (+12.0 \%) \\
W1 & Online Retail II H7 & absolute error & weekday median 8 & 0.787 /
0.763 & +3.1 \% {[}+1.3, +5.0{]} & 0.0056 & 6/6 & supported &
inconclusive (+4.9 \%) \\
\end{longtable}
}

\emph{Source: \texttt{results/analysis/analysis\_W3.json},
\texttt{analysis\_W1.json}; CSV:
\texttt{tables/tableA1\_replication\_H1.csv}.}

\textbf{Table A2. Seed sensitivity of the W2 item-day improvement (H1)}

{\def\LTcaptype{none} % do not increment counter
\begin{longtable}[]{@{}lllll@{}}
\toprule\noalign{}
cell & seed 42 (frozen primary) & seed 43 & seed 44 & range \\
\midrule\noalign{}
\endhead
\bottomrule\noalign{}
\endlastfoot
CA\_1 H28 & +4.58 \% & +4.38 \% & +4.57 \% & 0.20 pp \\
CA\_1 H7 & +3.95 \% & +4.25 \% & +4.20 \% & 0.30 pp \\
TX\_1 H28 & +2.53 \% & +2.37 \% & +2.30 \% & 0.23 pp \\
TX\_1 H7 & +4.54 \% & +4.53 \% & +4.41 \% & 0.12 pp \\
WI\_1 H28 & +3.33 \% & +3.53 \% & +3.73 \% & 0.40 pp \\
WI\_1 H7 & +6.02 \% & +6.07 \% & +6.11 \% & 0.09 pp \\
Online Retail II H28 & +5.81 \% & +5.62 \% & +5.95 \% & 0.33 pp \\
Online Retail II H7 & +4.02 \% & +4.17 \% & +3.81 \% & 0.37 pp \\
\end{longtable}
}

\emph{Source: \texttt{results/analysis/analysis\_W2.json}
(\texttt{seed\_sensitivity}); CSV: \texttt{tables/tableA2\_seeds.csv}.}

\textbf{Table A3 (exploratory). Segment results, window W2}

{\def\LTcaptype{none} % do not increment counter
\begin{longtable}[]{@{}
  >{\raggedright\arraybackslash}p{(\linewidth - 6\tabcolsep) * \real{0.2500}}
  >{\raggedright\arraybackslash}p{(\linewidth - 6\tabcolsep) * \real{0.2500}}
  >{\raggedright\arraybackslash}p{(\linewidth - 6\tabcolsep) * \real{0.2500}}
  >{\raggedright\arraybackslash}p{(\linewidth - 6\tabcolsep) * \real{0.2500}}@{}}
\toprule\noalign{}
\begin{minipage}[b]{\linewidth}\raggedright
cell
\end{minipage} & \begin{minipage}[b]{\linewidth}\raggedright
segment
\end{minipage} & \begin{minipage}[b]{\linewidth}\raggedright
series
\end{minipage} & \begin{minipage}[b]{\linewidth}\raggedright
item-day improvement vs matched baseline {[}95 \% CI{]}
\end{minipage} \\
\midrule\noalign{}
\endhead
\bottomrule\noalign{}
\endlastfoot
CA\_1 H28 & frequency=frequent & 46 & +8.1 \% {[}+3.5, +12.7{]} \\
CA\_1 H28 & frequency=intermittent\_or\_sparse & 304 & +2.5 \% {[}+0.5,
+4.7{]} \\
CA\_1 H28 & volume=high & 117 & +6.3 \% {[}+3.1, +9.7{]} \\
CA\_1 H28 & volume=low & 117 & +1.5 \% {[}−0.7, +4.0{]} \\
CA\_1 H28 & volume=mid & 116 & +1.0 \% {[}−1.1, +3.2{]} \\
CA\_1 H7 & frequency=frequent & 46 & +6.5 \% {[}+2.9, +11.1{]} \\
CA\_1 H7 & frequency=intermittent\_or\_sparse & 304 & +2.6 \% {[}+1.3,
+3.9{]} \\
CA\_1 H7 & volume=high & 117 & +5.0 \% {[}+2.7, +7.8{]} \\
CA\_1 H7 & volume=low & 117 & +2.9 \% {[}+1.2, +4.8{]} \\
CA\_1 H7 & volume=mid & 116 & +1.7 \% {[}+0.2, +3.1{]} \\
TX\_1 H28 & frequency=frequent & 33 & +4.4 \% {[}+0.7, +8.4{]} \\
TX\_1 H28 & frequency=intermittent\_or\_sparse & 317 & +1.8 \% {[}+0.3,
+3.3{]} \\
TX\_1 H28 & volume=high & 117 & +2.6 \% {[}+0.6, +4.8{]} \\
TX\_1 H28 & volume=low & 117 & +0.9 \% {[}−1.1, +3.1{]} \\
TX\_1 H28 & volume=mid & 116 & +2.9 \% {[}+1.2, +4.9{]} \\
TX\_1 H7 & frequency=frequent & 33 & +6.9 \% {[}+1.9, +12.4{]} \\
TX\_1 H7 & frequency=intermittent\_or\_sparse & 317 & +3.7 \% {[}+1.7,
+5.9{]} \\
TX\_1 H7 & volume=high & 117 & +5.8 \% {[}+3.2, +8.8{]} \\
TX\_1 H7 & volume=low & 117 & +0.7 \% {[}−1.7, +3.2{]} \\
TX\_1 H7 & volume=mid & 116 & +2.0 \% {[}+0.0, +4.2{]} \\
WI\_1 H28 & frequency=frequent & 41 & +4.9 \% {[}+0.4, +10.1{]} \\
WI\_1 H28 & frequency=intermittent\_or\_sparse & 308 & +2.6 \% {[}+1.0,
+4.3{]} \\
WI\_1 H28 & volume=high & 116 & +4.4 \% {[}+2.0, +7.1{]} \\
WI\_1 H28 & volume=low & 117 & +0.5 \% {[}−1.5, +2.8{]} \\
WI\_1 H28 & volume=mid & 116 & +1.3 \% {[}−0.9, +3.6{]} \\
WI\_1 H7 & frequency=frequent & 41 & +9.7 \% {[}+3.2, +18.0{]} \\
WI\_1 H7 & frequency=intermittent\_or\_sparse & 308 & +4.4 \% {[}+2.4,
+6.6{]} \\
WI\_1 H7 & volume=high & 116 & +8.5 \% {[}+4.8, +12.5{]} \\
WI\_1 H7 & volume=low & 117 & −0.5 \% {[}−2.5, +2.0{]} \\
WI\_1 H7 & volume=mid & 116 & +1.8 \% {[}−0.8, +4.6{]} \\
Online Retail II H28 & frequency=intermittent\_or\_sparse & 83 & +6.1 \%
{[}+2.4, +10.5{]} \\
Online Retail II H28 & volume=high & 33 & +5.9 \% {[}+2.5, +9.3{]} \\
Online Retail II H28 & volume=low & 34 & +5.9 \% {[}−3.0, +14.0{]} \\
Online Retail II H28 & volume=mid & 33 & +5.7 \% {[}−0.0, +12.4{]} \\
Online Retail II H7 & frequency=intermittent\_or\_sparse & 83 & +4.5 \%
{[}+1.8, +7.8{]} \\
Online Retail II H7 & volume=high & 33 & +3.7 \% {[}+1.0, +6.8{]} \\
Online Retail II H7 & volume=low & 34 & +5.0 \% {[}−0.3, +10.9{]} \\
Online Retail II H7 & volume=mid & 33 & +4.4 \% {[}+1.1, +9.5{]} \\
\end{longtable}
}

\emph{Source: \texttt{results/analysis/analysis\_W2.json}
(\texttt{segments\_exploratory}; B = 2,000, no multiplicity control);
CSV: \texttt{tables/tableA3\_segments.csv}.}

\textbf{Table A4. Design effect of the crossed series × week structure
(H1 comparison, W2)}

{\def\LTcaptype{none} % do not increment counter
\begin{longtable}[]{@{}
  >{\raggedright\arraybackslash}p{(\linewidth - 10\tabcolsep) * \real{0.1667}}
  >{\raggedright\arraybackslash}p{(\linewidth - 10\tabcolsep) * \real{0.1667}}
  >{\raggedright\arraybackslash}p{(\linewidth - 10\tabcolsep) * \real{0.1667}}
  >{\raggedright\arraybackslash}p{(\linewidth - 10\tabcolsep) * \real{0.1667}}
  >{\raggedright\arraybackslash}p{(\linewidth - 10\tabcolsep) * \real{0.1667}}
  >{\raggedright\arraybackslash}p{(\linewidth - 10\tabcolsep) * \real{0.1667}}@{}}
\toprule\noalign{}
\begin{minipage}[b]{\linewidth}\raggedright
cell
\end{minipage} & \begin{minipage}[b]{\linewidth}\raggedright
series
\end{minipage} & \begin{minipage}[b]{\linewidth}\raggedright
weekly blocks
\end{minipage} & \begin{minipage}[b]{\linewidth}\raggedright
series-week cells
\end{minipage} & \begin{minipage}[b]{\linewidth}\raggedright
design effect on the s.e.
\end{minipage} & \begin{minipage}[b]{\linewidth}\raggedright
approx. effective independent cells
\end{minipage} \\
\midrule\noalign{}
\endhead
\bottomrule\noalign{}
\endlastfoot
CA\_1 H28 & 350 & 24 & 8400 & 1.9 & 2321 \\
CA\_1 H7 & 350 & 24 & 8400 & 1.7 & 2800 \\
TX\_1 H28 & 350 & 24 & 8400 & 1.5 & 3668 \\
TX\_1 H7 & 350 & 24 & 8400 & 1.8 & 2558 \\
WI\_1 H28 & 349 & 24 & 8376 & 1.5 & 3980 \\
WI\_1 H7 & 349 & 24 & 8376 & 1.7 & 2856 \\
Online Retail II H28 & 100 & 24 & 2400 & 2.5 & 397 \\
Online Retail II H7 & 100 & 24 & 2400 & 2.0 & 610 \\
\end{longtable}
}

\emph{Source: \texttt{results/summary\_extra.json}
(\texttt{design\_effects\_W2}); CSV:
\texttt{tables/tableA4\_design\_effect.csv}.}

\subsection{References}\label{references}

\emph{Bibliographic metadata from Crossref (29 DOI entries;
\texttt{REFERENCE\_AUDIT.csv}); Ke et al.~checked on the NeurIPS
proceedings page, Damato et al.~on the arXiv API, Chen (2012) on the UCI
page; Koenker \& Bassett pages checked on the Econometric Society page
(2026-10-22). Claim-level support is in
\texttt{REFERENCE\_CLAIM\_CHECK.csv}: most claims were checked against
abstracts only; the M5 report was read in a pre-publication copy.}

Athanasopoulos, G., Hyndman, R.J., Kourentzes, N., \& Petropoulos, F.
(2017). Forecasting with temporal hierarchies. \emph{European Journal of
Operational Research}, 262, 60-74.
https://doi.org/10.1016/j.ejor.2017.02.046

Bojer, C.S., \& Meldgaard, J.P. (2021). Kaggle forecasting competitions:
An overlooked learning opportunity. \emph{International Journal of
Forecasting}, 37, 587-603.
https://doi.org/10.1016/j.ijforecast.2020.07.007

Cameron, A.C., Gelbach, J.B., \& Miller, D.L. (2011). Robust Inference
With Multiway Clustering. \emph{Journal of Business \& Economic
Statistics}, 29, 238-249. https://doi.org/10.1198/jbes.2010.07136

Chen, D. (2012). Online Retail II {[}Dataset{]}. UCI Machine Learning
Repository. https://doi.org/10.24432/C5CG6D (CC BY 4.0).

Chen, T., \& Guestrin, C. (2016). XGBoost: A scalable tree boosting
system. \emph{Proceedings of the 22nd ACM SIGKDD International
Conference on Knowledge Discovery and Data Mining}, 785-794.
https://doi.org/10.1145/2939672.2939785

Croston, J.D. (1972). Forecasting and Stock Control for Intermittent
Demands. \emph{Journal of the Operational Research Society}, 23,
289-303. https://doi.org/10.1057/jors.1972.50

Damato, S., Rubattu, N., Azzimonti, D., \& Corani, G. (2026).
Intermittent time series forecasting: local vs global models.
arXiv:2601.14031 (preprint; authors and title per the arXiv API,
2026-10-22).

Diebold, F.X., \& Mariano, R.S. (1995). Comparing Predictive Accuracy.
\emph{Journal of Business \& Economic Statistics}, 13, 253-263.
https://doi.org/10.1080/07350015.1995.10524599

Fildes, R., Ma, S., \& Kolassa, S. (2022). Retail forecasting: Research
and practice. \emph{International Journal of Forecasting}, 38,
1283-1318. https://doi.org/10.1016/j.ijforecast.2019.06.004

Friedman, J.H. (2001). Greedy function approximation: A gradient
boosting machine. \emph{The Annals of Statistics}, 29, 1189-1232.
https://doi.org/10.1214/aos/1013203451

Gneiting, T. (2011). Making and Evaluating Point Forecasts.
\emph{Journal of the American Statistical Association}, 106, 746-762.
https://doi.org/10.1198/jasa.2011.r10138

Hewamalage, H., Ackermann, K., \& Bergmeir, C. (2023). Forecast
evaluation for data scientists: common pitfalls and best practices.
\emph{Data Mining and Knowledge Discovery}, 37, 788-832.
https://doi.org/10.1007/s10618-022-00894-5

Hyndman, R.J., \& Koehler, A.B. (2006). Another look at measures of
forecast accuracy. \emph{International Journal of Forecasting}, 22,
679-688. https://doi.org/10.1016/j.ijforecast.2006.03.001

Hyndman, R.J., Ahmed, R.A., Athanasopoulos, G., \& Shang, H.L. (2011).
Optimal combination forecasts for hierarchical time series.
\emph{Computational Statistics \& Data Analysis}, 55, 2579-2589.
https://doi.org/10.1016/j.csda.2011.03.006

Januschowski, T., Wang, Y., Torkkola, K., Erkkilä, T., Hasson, H., \&
Gasthaus, J. (2022). Forecasting with trees. \emph{International Journal
of Forecasting}, 38, 1473-1481.
https://doi.org/10.1016/j.ijforecast.2021.10.004

Ke, G., Meng, Q., Finley, T., Wang, T., Chen, W., Ma, W., Ye, Q., \&
Liu, T.-Y. (2017). LightGBM: A highly efficient gradient boosting
decision tree. \emph{Advances in Neural Information Processing Systems},
30 (authors and title checked on the NeurIPS proceedings page,
2026-10-22).

Koenker, R., \& Bassett, G. (1978). Regression Quantiles.
\emph{Econometrica}, 46, 33-50. https://doi.org/10.2307/1913643

Kolassa, S. (2016). Evaluating predictive count data distributions in
retail sales forecasting. \emph{International Journal of Forecasting},
32, 788-803. https://doi.org/10.1016/j.ijforecast.2015.12.004

Kolassa, S. (2020). Why the ``best'' point forecast depends on the error
or accuracy measure. \emph{International Journal of Forecasting}, 36,
208-211. https://doi.org/10.1016/j.ijforecast.2019.02.017

Kourentzes, N. (2014). On intermittent demand model optimisation and
selection. \emph{International Journal of Production Economics}, 156,
180-190. https://doi.org/10.1016/j.ijpe.2014.06.007

Makridakis, S., Spiliotis, E., \& Assimakopoulos, V. (2018). Statistical
and Machine Learning forecasting methods: Concerns and ways forward.
\emph{PLOS ONE}, 13, e0194889.
https://doi.org/10.1371/journal.pone.0194889

Makridakis, S., Spiliotis, E., \& Assimakopoulos, V. (2020). The M4
Competition: 100,000 time series and 61 forecasting methods.
\emph{International Journal of Forecasting}, 36, 54-74.
https://doi.org/10.1016/j.ijforecast.2019.04.014

Makridakis, S., Spiliotis, E., \& Assimakopoulos, V. (2022). M5 accuracy
competition: Results, findings, and conclusions. \emph{International
Journal of Forecasting}, 38, 1346-1364.
https://doi.org/10.1016/j.ijforecast.2021.11.013

Nikolopoulos, K., Syntetos, A.A., Boylan, J.E., Petropoulos, F., \&
Assimakopoulos, V. (2011). An aggregate--disaggregate intermittent
demand approach (ADIDA) to forecasting: an empirical proposition and
analysis. \emph{Journal of the Operational Research Society}, 62,
544-554. https://doi.org/10.1057/jors.2010.32

Owen, A.B. (2007). The pigeonhole bootstrap. \emph{The Annals of Applied
Statistics}, 1, 386-411. https://doi.org/10.1214/07-AOAS122

Owen, A.B., \& Eckles, D. (2012). Bootstrapping data arrays of arbitrary
order. \emph{The Annals of Applied Statistics}, 6, 895-927.
https://doi.org/10.1214/12-AOAS547

Politis, D.N., \& Romano, J.P. (1994). The Stationary Bootstrap.
\emph{Journal of the American Statistical Association}, 89, 1303-1313.
https://doi.org/10.1080/01621459.1994.10476870

Prestwich, S., Rossi, R., Armagan Tarim, S., \& Hnich, B. (2014).
Mean-based error measures for intermittent demand forecasting.
\emph{International Journal of Production Research}, 52, 6782-6791.
https://doi.org/10.1080/00207543.2014.917771

Spiliotis, E., Abolghasemi, M., Hyndman, R.J., Petropoulos, F., \&
Assimakopoulos, V. (2021). Hierarchical forecast reconciliation with
machine learning. \emph{Applied Soft Computing}, 112, 107756.
https://doi.org/10.1016/j.asoc.2021.107756

Spiliotis, E., Makridakis, S., Semenoglou, A.A., \& Assimakopoulos, V.
(2022). Comparison of statistical and machine learning methods for daily
SKU demand forecasting. \emph{Operational Research}, 22, 3037-3061.
https://doi.org/10.1007/s12351-020-00605-2

Syntetos, A.A., \& Boylan, J.E. (2005). The accuracy of intermittent
demand estimates. \emph{International Journal of Forecasting}, 21,
303-314. https://doi.org/10.1016/j.ijforecast.2004.10.001

Teunter, R.H., Syntetos, A.A., \& Zied Babai, M. (2011). Intermittent
demand: Linking forecasting to inventory obsolescence. \emph{European
Journal of Operational Research}, 214, 606-615.
https://doi.org/10.1016/j.ejor.2011.05.018

\end{document}
